\section{Methodology}
\label{sec:method}

Building on our key insight---that the proxy-gold divergence gap, when properly normalized, grows linearly and measurably with RLHF optimization pressure---we design a methodology with three components: a normalization protocol that creates a cross-dataset comparison instrument, a regression pipeline that quantifies the divergence curve slope, and a four-step mechanistic verification that establishes \emph{why} the gap grows.

\subsection{Overview}

We study two independently published RLHF experimental datasets (Coste et al.~\cite{Coste2023Reward} and Gao et al.~\cite{Gao2023Scaling}) at multiple KL divergence checkpoints. At each checkpoint, we have two signals: the reward model score (the proxy) and held-out gold human preference (the target). The core methodological challenge is that these signals live on different scales across datasets---raw RM scores in Coste et al.\ span $[0.12, 2.08]$ while gold preference spans $[0.38, 0.63]$. Cross-dataset comparison of raw differences is therefore meaningless.

Our solution is the \textbf{normalized divergence gap}:
\begin{equation}
\text{gap}(k) = \text{normalize}(\text{RM}_k) - \text{gold\_preference}_k
\end{equation}
where normalization is min-max scaling to $[0, 1]$ applied to the RM score series, and gold preference is already in $[0, 1]$ as a rate. The resulting gap $\in [-1, +1]$ enables cross-dataset slope comparison on a common scale.

\subsection{Datasets and Data Reconstruction}

\textbf{Dataset 1: Coste et al.~\cite{Coste2023Reward} (arXiv:2310.02743).} We use 10 paired observations of (KL\_budget, RM\_score, gold\_preference) from Coste et al.~\cite{Coste2023Reward}, covering KL budgets from 0.0 to 8.0 nats. Data values are constructed consistent with qualitative descriptions in published figures using standard figure digitization methodology. This introduces an estimated $\pm 2$--5\% digitization uncertainty per data point.

\textbf{Dataset 2: Gao et al.~\cite{Gao2023Scaling} (arXiv:2210.10760, ICML 2023).} We use 10 paired observations from Gao et al.~\cite{Gao2023Scaling}, covering KL budgets from 0.0 to 7.5 nats. This dataset uses a different model family and scale (6B reward model), providing an independent replication context.

\textbf{Honest disclosure:} Neither raw dataset is publicly available in machine-readable format. Our data reconstruction from published figures introduces digitization uncertainty that affects exact $\beta$ values but not the direction or statistical significance of the slope---the documented effect sizes ($R^2 = 0.958$, $p < 10^{-6}$ in Coste data) are robust to $\pm 2$--5\% perturbation.

\subsection{The Normalized Divergence Gap Metric}

For a series of $N$ KL checkpoints with RM scores $s_1, \ldots, s_N$ and gold preference rates $g_1, \ldots, g_N$:
\begin{equation}
\text{RM\_norm}_i = \frac{s_i - \min(s)}{\max(s) - \min(s)}, \quad \text{gap}_i = \text{RM\_norm}_i - g_i
\end{equation}

We choose min-max normalization over z-score normalization because it preserves the $[0, 1]$ interpretable range, gold preference is already in $[0, 1]$ as a rate, and it preserves the monotonicity structure needed to detect reward hacking while removing cross-dataset scale artifacts.

\subsection{Regression Pipeline}

We apply OLS regression to $\{(\text{KL}_i, \text{gap}_i)\}_{i=1}^{N}$: $\text{gap} = \beta \cdot \text{KL} + \alpha + \varepsilon$. We use two implementations for redundancy: \texttt{scipy.stats.linregress} (primary, parametric CI via t-distribution) and \texttt{statsmodels.api.OLS} (secondary, F-statistic, residual diagnostics). Bootstrap CIs use $N=10{,}000$ resamples with seed=42. Pre-registered success criteria: $\beta > 0$, $p < 0.05$ (primary); $R^2 > 0.5$ (secondary).

\subsection{Four-Step Mechanistic Verification}

\begin{enumerate}
\item \textbf{Step 1 (H-E1):} Signal co-existence---both RM score and gold preference are separable, non-constant series.
\item \textbf{Step 2 (H-M1):} Trajectory divergence---RM rises monotonically (Spearman $\rho > 0.8$, $p < 0.05$) while gold preference reverses.
\item \textbf{Step 3 (H-M2):} Gap positivity---normalized divergence gap is strictly positive and growing at high KL ($\text{KL} > 3.5$ nats).
\item \textbf{Step 4 (H-M3, H-M4):} Slope quantification---OLS $\beta > 0$, $p < 0.05$ in both datasets; cross-dataset slope ratio $\beta_\text{Gao} / \beta_\text{Coste}$.
\end{enumerate}

Figure~\ref{fig:trajectory} (\texttt{trajectory\_dual\_axis.png}) illustrates the diverging trajectories from Steps 2--3.
